\documentclass[../Main.tex]{subfiles}

\begin{document}
In this chapter we will require normalised wavefunctions unless otherwise specified.
\section{Bound States}
\subsection{Infinite Potential Well}
Consider a potential of the form:
\begin{equation*}
    U(x) =
    \begin{cases}
        0 & \abs{x} \leq a \\
        \infty & \abs{x} > a
    \end{cases}
\end{equation*}
which means that for $\abs{x} > a$ we need $\chi(x) = 0$. The solutions are:
\begin{equation*}
    \chi_n(x) =
    \begin{cases}
        a^{-\frac12} \sin\left(\frac{n\pi x}{2a}\right) & \text{$n$ even} \\
        a^{-\frac12} \cos\left(\frac{n\pi x}{2a}\right) & \text{$n$ odd}
    \end{cases}
\end{equation*}
found by solving the second-order TDSE and applying boundary conditions. Note that the ground state is at $n = 1$ not $n = 0$ because this is the first normalisable wavefunction.
\begin{proposition}
    For a quantum system with non-degenerate eigenfunctions and even potential ($U(x) = U(-x)$), the eigenfunctions of $\hat{H}$ must be either odd or even.
    \label{propEvenOddWF}
\end{proposition}
\subsection{Finite Potential Well}
In this case we have finite potential and so we require $\chi \in C^1(\R)$.
\begin{equation*}
    U(x) =
    \begin{cases}
        0 & \abs{x} \leq a \\
        U_0 & \abs{x} > a
    \end{cases}
\end{equation*}
We deal with bound states with $E < U_0$. We have even and odd state for $\chi$. The even states are:
\begin{equation*}
    \chi(x) =
    \begin{cases}
        Ce^{\bar{k}x} & x < -a \\
        B\cos(kx) & \abs{x} < a \\
        Ce^{-\bar{k}x} & x > a
    \end{cases}
\end{equation*}
where $k\tan(ka) = \bar{k}$, $k^2 + \bar{k}^2 = \frac{2mU_0}{\hbar^2}$. This has finitely many solutions for each $\bar{k}$. $A, B$ are determined by normalisation and the condition that the pieces must agree at $x = a$.

The odd states are:
\begin{equation*}
    \chi(x) =
    \begin{cases}
        Ce^{-\bar{k}x} & x < -a \\
        B\sin(kx) & \abs{x} < a \\
        -Ce^{\bar{k}x} & x > a
    \end{cases}
\end{equation*}
where $k = -\bar{k} \tan(ka)$ and once again $A$, $B$ are normalisation constants. In this case we have no solutions for $\frac{8mU_0}{\hbar^2 \pi^2} < a^{-2}$.
\subsection{Quantum Harmonic Oscillator}
Now we take:
\begin{equation*}
    U(x) = \frac12kx^2
\end{equation*}
and we solve with a power series (and require that it terminates) to get:
\begin{equation*}
    \chi_n(\xi) = f_n(\xi) e^{-\xi^2 / 2}
\end{equation*}
where $\xi = \sqrt{m\omega}{\hbar}$ and $f_n$ are the Hermite polynomials satisfying:
\begin{equation*}
    f_n(\xi) = (-1)^n e^{\xi^2} \frac{d^{n}}{d\xi^{n}}\left(e^{-\xi^2}\right)
\end{equation*}
and the eigenvalues are $E_n = \left(n + \frac12\right)\hbar\omega$.
\section{Free Particle}
The solution to the TISE for a free particle is $\chi_k(x) = e^{ikx}$. This is non-normalisable and so we can either:
\begin{enumerate}
    \item Build a superposition of particles that is normalisable
    \item Change the interpretation
\end{enumerate}
\subsection{Gaussian Wavepacket}
The Gaussian wavepacket is a linear superposition of solutions to the TISE with the following form:
\begin{equation*}
    \psi(x, t) = \int_{-\infty}^{\infty} A_{GP}(k) \psi_k(x, t) dk
\end{equation*}
where:
\begin{equation*}
    A_{GP}(k) = \exp\left[-\frac{\sigma}{2}(k - k_0)^2\right],
\end{equation*}
a normal distribution $N(k_0, \sigma^2)$. The wavefunction with this form is now normalisable.

The expected position is $\E{\hat{x}}_{GP} = \frac{\hbar k_0}{m}t$, with uncertainty: 
\begin{equation*}
    \Delta x = \sqrt{\frac12 \left(\sigma + \frac{\hbar^2 t^2}{m^2 \sigma}\right)}
\end{equation*}
which spreads out over time.

Expected momentum is $\E{\hat{p}}_{GP} = \hbar k_0$, with uncertainty $\frac{\hbar}{\sqrt{2\sigma}}$. In fact, Heisenberg Uncertainty Principle bound is attained in this case:
\begin{equation*}
    \Delta x \Delta p \geq \frac{\hbar}{2}.
\end{equation*}
\subsection{Beam Interpretation}
We can instead change our interpretation of the wavefunction $\chi_k(x) = e^{ikx}$. We say that this is a beam of particles with average density $\abs{\psi}^2$. The full wavefunction is 
\begin{equation*}
    \psi_k(x, t) = Ae^{ikx} e^{-i\hbar k^2 / 2m}.
\end{equation*}
Probability current is given by:
\begin{align*}
    j_k(x, t) &= -\frac{i\hbar}{2m} \left(\psi_k^* \frac{\partial \psi_k}{\partial x} - \psi_k \frac{\partial \psi_k^*}{\partial x}\right) \\
    &= \abs{A}^2 \frac{\hbar k}{m} \\
    &= \abs{A}^2 \frac{p_k}{m}
\end{align*}
and this is understood as the average flux of particles (average density times velocity).
\section{Scattering States}
\begin{definition}{Scattering state}
    A \underline{scattering state} is a quantum system with a wavefunction that does not decay to $0$ as $\abs{x} \to \infty$.
\end{definition}
\subsection{Potential Step}
Consider:
\begin{equation*}
    U(x) =
    \begin{cases}
        0 & x \leq 0 \\
        U_0 & x > 0
    \end{cases}
\end{equation*}
Then solving on the two regions for a beam incident from $-\infty$,
\begin{equation*}
    \chi_{k, \bar{k}}(x) =
    \begin{cases}
        Ae^{ikx} + Be^{-ikx} & x \leq 0 \\
        Ce^{i\bar{k}x} & x > 0
    \end{cases}
\end{equation*}
where $B, C$ can be determined in terms of $A$ using continuity of $\chi$ and its derivative at zero. $k$ and $\bar{k}$ depend on the energy:
\begin{equation*}
    k = \sqrt{\frac{2mE}{\hbar^2}},\qquad \bar{k} = \sqrt{\frac{2m(E - U_0)}{\hbar^2}}
\end{equation*}
The flux of particles can be split into $j_{inc}, j_{ref}$ and $j_{trn}$ for the incident, reflected and transmitted beams.
\begin{align*}
    j_{inc} &= \frac{\hbar k}{m} \abs{A}^2 \\
    j_{ref} &= \frac{\hbar k}{m} \abs{B}^2 \\
    j_{trn} &= \frac{\hbar k}{m} \abs{C}^2 \\
\end{align*}
and we define the reflection and transmission coefficients:
\begin{align*}
    R &= \frac{j_{ref}}{j_{inc}} \\
    T &= \frac{j_{trn}}{j_{inc}}
\end{align*}
which satisfy $R + T = 1$.

In the case $E < U_0$ we find a similar solution but with $0$ momentum on $x > 0$ and a transmission coefficient of $0$.
\subsection{Potential Barrier}
Now consider a potential:
\begin{equation*}
    U(x) =
    \begin{cases}
        U_0 & x \in (0, a) \\
        0 & \text{otherwise}
    \end{cases}
\end{equation*}
and define:
\begin{equation*}
    k = \sqrt{\frac{2mE}{\hbar^2}},\qquad \eta = \sqrt{\frac{2m(U_0 - E)}{\hbar^2}}
\end{equation*}
We find the solution:
\begin{equation*}
    \chi(x) =
    \begin{cases}
        e^{ikx} + Ae^{-ikx} & x \leq 0 \\
        Be^{-\eta x} + Ce^{\eta x} & x \in (0, a) \\
        De^{ikx} & x \geq a
    \end{cases}
\end{equation*}
where $A, B, C$ and $D$ can be determined using continuity of $\chi$ and $\chi'$ at the joins. The transmission coefficient is $\abs{D}^2$:
\begin{equation*}
    T = \frac{4k^2\eta^2}{(k^2+\eta^2)\sinh(\eta a) + 4k^2\eta^2}
\end{equation*}
which is non-zero, so particles can quantum tunnel through the barrier despite having less energy.
\end{document}